% TOP_PROBLEM: {"id": 1216, "title": "Kalai's $3^d$ Conjecture", "category_id": 6, "category": {"id": 6, "name": "geometry", "display_name": "Geometry", "description": "Euclidean and non-Euclidean geometry, geometric structures.", "slug": "geometry", "order_index": 6, "created_at": "2026-07-31T15:26:25.671Z"}, "status": "open"}
% FIRST_POSTED: 2026-10-01
% ATTEMPT_STATUS: unresolved
% Imported from: unsolvedmath_additional_500.tex; UnsolvedMath ID 1216
% !TeX program = lualatex
% Compile twice: lualatex 1216.tex
% Self-contained: no external bibliography, images, or \input files.
% Numeric section labels and visible numbers are original UnsolvedMath IDs.
\documentclass[11pt,a4paper]{article}
\usepackage[margin=24mm,headheight=14pt]{geometry}
\usepackage{fontspec}
\setmainfont{Latin Modern Roman}
\setsansfont{Latin Modern Sans}
\setmonofont{Latin Modern Mono}
\usepackage{amsmath,amssymb,mathtools}
\usepackage{unicode-math}
\setmathfont{Latin Modern Math}
\usepackage{microtype}
\usepackage{enumitem,xurl,needspace}
\usepackage[dvipsnames]{xcolor}
\usepackage{hyperref}
\hypersetup{unicode=true,colorlinks=true,linkcolor=MidnightBlue,urlcolor=MidnightBlue,
 pdftitle={UnsolvedMath Problem 1216},pdfauthor={Source-annotated compilation},bookmarksnumbered=true}
\usepackage{bookmark,tocloft,titlesec,fancyhdr}
\setlength{\cftsecnumwidth}{4.2em}
\cftsetpnumwidth{2.5em}
\cftsetrmarg{4em}
\titleformat{\section}{\Large\bfseries\raggedright}{\thesection}{0.7em}{}
\pagestyle{fancy}\fancyhf{}
\fancyhead[L]{\small\sffamily Additional UnsolvedMath statements}
\fancyhead[R]{\small Original catalogue IDs}
\fancyfoot[C]{\thepage}
\setlength{\parindent}{0pt}
\setlength{\parskip}{5pt plus 1pt}
\setlength{\emergencystretch}{3em}
\setcounter{tocdepth}{1}\setcounter{secnumdepth}{1}
\setlist[itemize]{leftmargin=3em,itemsep=4pt,topsep=4pt,parsep=0pt}
\allowdisplaybreaks
\newcommand{\N}{\mathbb{N}}
\newcommand{\Z}{\mathbb{Z}}
\newcommand{\Q}{\mathbb{Q}}
\newcommand{\R}{\mathbb{R}}
\newcommand{\C}{\mathbb{C}}
\newcommand{\F}{\mathbb{F}}
\newcommand{\A}{\mathbb{A}}
\newcommand{\PP}{\mathbb{P}}
\newcommand{\E}{\mathbb{E}}
\newcommand{\eps}{\varepsilon}
\newcommand{\dd}{\,\mathrm d}
\newcommand{\sref}[2]{\hyperlink{source:#1:#2}{[S#2]}}
\newif\ifshowcatalogue
\showcataloguetrue
\newcommand{\cataloguescope}[1]{\ifshowcatalogue
 \paragraph{Statement recorded in the supplied source.}
 #1\fi}
\begin{document}
% UnsolvedMath ID 1216; source catalogue code GEO-025

\setcounter{section}{1215}

\section{Kalai's 3\textasciicircum{}d Face Conjecture}\label{1216}

{\small\sffamily UnsolvedMath ID 1216 · Geometry}

\subsection{Definitions and mathematical statement}

\paragraph{Shared notation.}Unless a section says otherwise, $\N=\{1,2,\ldots\}$,
$\Z$ is the ring of integers, $\Q,\R,\C$ are the rational, real, and complex
fields, $[N]=\{1,\ldots,\lfloor N\rfloor\}$, and $\log$ is the natural
logarithm. For real functions of a parameter tending to infinity,
$f\ll g$ and $f=O(g)$ mean $|f|\leq Cg$ eventually for some fixed $C$;
$f\gg g$ reverses this comparison; $f=o(g)$ means $f/g\to0$;
$f\sim g$ means $f/g\to1$; and $f\asymp g$ means both order bounds.
A subscript on $O$ or $\ll$ specifies allowed constant dependence.
The expression $n^{o(1)}$ in an upper bound means growth smaller than
$n^\epsilon$ for each fixed $\epsilon>0$. Local definitions govern any
exception, finite-field convention, or use of the same symbol for another
object. An asymptotic claim is not automatically an effective algorithm.



A $d$-dimensional convex polytope is centrally symmetric if, after translating its centre to the origin, $P=-P$. Faces are intersections with supporting hyperplanes, together with the polytope itself. Count all nonempty faces, of every dimension from 0 to $d$, but not the empty face. The proposed bound is $\sum_{j=0}^{d}f_j(P)\geq3^d$. \sref{1216}{1} \sref{1216}{2}

\paragraph{Statement recorded in the supplied source.}
Does every centrally symmetric $d$-dimensional polytope have at least $3^d$ faces? \sref{1216}{1}

\subsection{Short English statement}

Does every centrally symmetric polytope have at least three to the power of its dimension nonempty faces in total?

\subsection{Research attempt: equality under products and polarity, and the first three dimensions}
\textbf{Status and source audit.} This proves exact equality for the product/polar family generated by a segment, and the lower bound through dimension three. It does not establish the general inequality. The inspected abstract of [S2] assumes symmetry in $d$ mutually orthogonal hyperplanes; those additional symmetries are not consequences of central symmetry. Its abstract also uses the word ``facets'' where the catalogue concerns all nonempty faces; a cube has only $2d$ facets, so that word cannot be adopted as the statement counted here.

\subsubsection{1. A multiplicative face invariant}
Write $F(P)$ for the number of nonempty faces including $P$. A supporting linear functional on $P\times Q$ splits as $\ell(x,y)=u(x)+v(y)$, and its maximizing face is $P^u\times Q^v$. Conversely every pair of nonempty faces is obtained this way, including the whole factor by taking the zero functional. Distinct pairs give distinct products by projection. Thus
\[F(P\times Q)=F(P)F(Q).\tag{1}\]
For a polytope containing 0 in its interior, polarity reverses the proper nonempty face lattice. It exchanges the empty face and the whole polytope in the augmented lattice. Thus proper nonempty faces correspond bijectively and each polytope contributes one whole face, giving
\[F(P^\circ)=F(P).\tag{2}\]
This is the standard supporting-face duality: a face corresponds to the functionals in $P^\circ$ taking value 1 throughout it, and applying the same operation again recovers the original face.

\subsubsection{2. Equality survives two nontrivial operations}
A segment has three nonempty faces. Starting from centrally symmetric segments and repeatedly taking Cartesian products and polar duals, central symmetry is preserved, dimensions add under products and are unchanged under polarity. Equations (1)--(2) prove by induction that every polytope so obtained satisfies $F(P)=3^{\dim P}$. In particular the cube has
\[f_j=2^{d-j}\binom dj\quad(0\leq j\leq d),\qquad
F(P)=\sum_j2^{d-j}\binom dj=3^d,\]
and its polar cross-polytope has the same total. This explains why both apparently opposite face distributions attain the proposed minimum.

\subsubsection{3. All centrally symmetric polytopes in dimensions at most three}
Dimension one is the segment. In dimension two central symmetry pairs vertices, and a full-dimensional polygon has at least four. For $v$ vertices it has $v$ edges and one whole face, hence $F(P)=2v+1\geq9$.

In dimension three write $v,e,f$ for the numbers of vertices, edges and facets. Euler's polyhedron formula $v-e+f=2$ gives $F(P)=2e+3$. This standard topological formula is the only external input in this low-dimensional count. Vertices come in antipodal pairs, and at least three pairs are needed to span three-dimensional space. If $v=6$, the vertices are $\pm a,\pm b,\pm c$ with $a,b,c$ linearly independent, so $P$ is a linear image of the octahedron and has $e=12$. If $v\geq8$, every vertex has degree at least three, so $2e\geq3v\geq24$. The minimum-degree assertion follows because the vertex figure is a two-dimensional polygon, having at least three vertices, one for each incident edge. In either case $F(P)\geq27$.

\subsubsection{4. Failure of the naive induction}
Not every centrally symmetric polytope is a product or a polar of a previously treated product. Nor do its facets have to be centrally symmetric: an octahedron already has triangular facets. Applying the dimension-$d-1$ assertion separately to facets would therefore use a false hypothesis and would also overcount their intersections. The diagnostic checks the displayed cube and cross-polytope face formulas through dimension 12 and the product identity on their products. No classification of all equality cases is claimed. The remaining task is a face-count mechanism valid for arbitrary centrally symmetric polytopes that neither assumes decomposition nor assumes symmetry of every proper face.

\subsection{Sources}

{\small

\begin{itemize}

\item[\hypertarget{source:1216:1}{S1}] ULAM.AI, \emph{UnsolvedMath}, supplied \texttt{problems.json}, record 1216 (GEO-025), statement and background. The embedded literature-review date is 2026-08-17; that is the archive’s review date, not a fresh certification. Dataset landing page: \url{https://huggingface.co/datasets/ulamai/UnsolvedMath}.

\item[\hypertarget{source:1216:2}{S2}] Gregory Chambers and Elia Portnoy, A note on Kalai's 3\textasciicircum{}d Conjecture, arXiv:2211.09215 (2022). \url{https://arxiv.org/abs/2211.09215}. Bibliographic information imported from the supplied record.

\item[\hypertarget{source:1216:3}{S3}] Isabella Novik, A tale of centrally symmetric polytopes and spheres, arXiv:1711.09310. \url{https://arxiv.org/abs/1711.09310}. Bibliographic information imported from the supplied record.

\end{itemize}

}

\label{end:1216}
\end{document}
